\small
\begin{enumerate}[label=\textbf{\arabic*.}, leftmargin=0pt, itemsep=1.2em]
\item Что в Прологе называется фактом?
\begin{itemize}[label=$\square$, leftmargin=1.5cm, itemsep=0.2em, topsep=0.3em]
\item Клауза с головой и непустым телом, истинная при выполнении условий
\item Клауза без тела, утверждение, истинное безусловно (например, parent(tom, bob).)
\item Вопрос к базе знаний, вводимый после приглашения \texttt{?-}
\item Директива, выполняемая при загрузке файла
\end{itemize}
{\small\textit{Правильный ответ: Клауза без тела, утверждение, истинное безусловно (например, parent(tom, bob).).}}

\item Какая последовательность символов отделяет голову правила от его тела?
\begin{itemize}[label=$\square$, leftmargin=1.5cm, itemsep=0.2em, topsep=0.3em]
\item :-
\item ->
\item :=
\item <-
\end{itemize}
{\small\textit{Правильный ответ: :-.}}

\item Какой запрос загружает файл family.pl в интерактивной среде SWI-Prolog?
\begin{itemize}[label=$\square$, leftmargin=1.5cm, itemsep=0.2em, topsep=0.3em]
\item \texttt{?-} [family].
\item \texttt{?-} load family.
\item \texttt{?-} open(family).
\item \texttt{?-} import(family).
\end{itemize}
{\small\textit{Правильный ответ: \texttt{?-} [family].}}

\item Каким встроенным предикатом завершают работу интерпретатора SWI-Prolog?
\begin{itemize}[label=$\square$, leftmargin=1.5cm, itemsep=0.2em, topsep=0.3em]
\item exit.
\item quit.
\item halt.
\item stop.
\end{itemize}
{\small\textit{Правильный ответ: halt.}}

\item Что ответит SWI-Prolog на запрос \texttt{?-} parent(bob, tom). для программы family.pl?
\begin{itemize}[label=$\square$, leftmargin=1.5cm, itemsep=0.2em, topsep=0.3em]
\item true.
\item false.
\item ERROR: Unknown procedure
\item X = bob.
\end{itemize}
{\small\textit{Правильный ответ: false.}}

\item Какую клавишу нажимают в top-level SWI-Prolog, чтобы получить следующее решение запроса?
\begin{itemize}[label=$\square$, leftmargin=1.5cm, itemsep=0.2em, topsep=0.3em]
\item ; (точка с запятой)
\item Enter
\item . (точка)
\item Esc
\end{itemize}
{\small\textit{Правильный ответ: ; (точка с запятой).}}

\item Как называется механизм возврата к последней точке выбора для поиска альтернативного решения?
\begin{itemize}[label=$\square$, leftmargin=1.5cm, itemsep=0.2em, topsep=0.3em]
\item Унификация
\item Хвостовая рекурсия
\item Возврат (backtracking)
\item Индексация клауз
\end{itemize}
{\small\textit{Правильный ответ: Возврат (backtracking).}}

\item Что такое унификация в Прологе?
\begin{itemize}[label=$\square$, leftmargin=1.5cm, itemsep=0.2em, topsep=0.3em]
\item Объединение двух файлов программы в один
\item Сопоставление двух термов с подбором значений переменных, при которых термы становятся одинаковыми
\item Приведение всех атомов к нижнему регистру
\item Удаление повторяющихся фактов из базы
\end{itemize}
{\small\textit{Правильный ответ: Сопоставление двух термов с подбором значений переменных, при которых термы становятся одинаковыми.}}

\item Какие имена в Прологе обозначают переменные?
\begin{itemize}[label=$\square$, leftmargin=1.5cm, itemsep=0.2em, topsep=0.3em]
\item Имена, начинающиеся со строчной буквы
\item Имена, начинающиеся с заглавной буквы или символа подчёркивания
\item Имена, заключённые в одинарные кавычки
\item Имена, начинающиеся с цифры
\end{itemize}
{\small\textit{Правильный ответ: Имена, начинающиеся с заглавной буквы или символа подчёркивания.}}

\item Чем является tom в факте parent(tom, bob).?
\begin{itemize}[label=$\square$, leftmargin=1.5cm, itemsep=0.2em, topsep=0.3em]
\item Переменной
\item Строкой
\item Атомом (константой)
\item Именем предиката
\end{itemize}
{\small\textit{Правильный ответ: Атомом (константой).}}

\item Какой запрос включает пошаговую трассировку выполнения целей?
\begin{itemize}[label=$\square$, leftmargin=1.5cm, itemsep=0.2em, topsep=0.3em]
\item \texttt{?-} step.
\item \texttt{?-} trace.
\item \texttt{?-} debug\_on.
\item \texttt{?-} verbose.
\end{itemize}
{\small\textit{Правильный ответ: \texttt{?-} trace.}}

\item Какие порты процедурной модели (модели Бёрда) показывает трассировщик SWI-Prolog?
\begin{itemize}[label=$\square$, leftmargin=1.5cm, itemsep=0.2em, topsep=0.3em]
\item Call, Exit, Redo, Fail
\item Start, Stop, Pause, Resume
\item Input, Output, Error, Return
\item Load, Run, Trace, Halt
\end{itemize}
{\small\textit{Правильный ответ: Call, Exit, Redo, Fail.}}

\item Какую роль играет первая клауза ancestor(X, Y) :- parent(X, Y). в определении ancestor/2?
\begin{itemize}[label=$\square$, leftmargin=1.5cm, itemsep=0.2em, topsep=0.3em]
\item Задаёт рекурсивный шаг
\item Задаёт базовый случай рекурсии: родитель является предком
\item Запрещает циклы в родословной
\item Объявляет типы аргументов
\end{itemize}
{\small\textit{Правильный ответ: Задаёт базовый случай рекурсии: родитель является предком.}}

\item Какое первое решение вернёт запрос \texttt{?-} grandparent(X, jim).?
\begin{itemize}[label=$\square$, leftmargin=1.5cm, itemsep=0.2em, topsep=0.3em]
\item X = tom
\item X = bob
\item X = pat
\item false.
\end{itemize}
{\small\textit{Правильный ответ: X = bob.}}

\item Сколько решений имеет запрос \texttt{?-} parent(X, Y). для программы family.pl?
\begin{itemize}[label=$\square$, leftmargin=1.5cm, itemsep=0.2em, topsep=0.3em]
\item 3
\item 4
\item 5
\item 6
\end{itemize}
{\small\textit{Правильный ответ: 5.}}

\item Что означает ответ false. после нескольких нажатий ; в top-level?
\begin{itemize}[label=$\square$, leftmargin=1.5cm, itemsep=0.2em, topsep=0.3em]
\item В программе есть ошибка
\item Найденные ранее решения неверны
\item Других решений нет, перебор исчерпан
\item Запрос содержит синтаксическую ошибку
\end{itemize}
{\small\textit{Правильный ответ: Других решений нет, перебор исчерпан.}}

\item Что делает запрос \texttt{?-} listing(ancestor/2).?
\begin{itemize}[label=$\square$, leftmargin=1.5cm, itemsep=0.2em, topsep=0.3em]
\item Удаляет предикат ancestor/2 из базы
\item Выводит на экран все клаузы предиката ancestor/2
\item Показывает количество решений ancestor/2
\item Включает трассировку ancestor/2
\end{itemize}
{\small\textit{Правильный ответ: Выводит на экран все клаузы предиката ancestor/2.}}

\item Для чего служит запрос \texttt{?-} make.?
\begin{itemize}[label=$\square$, leftmargin=1.5cm, itemsep=0.2em, topsep=0.3em]
\item Компилирует программу в исполняемый файл
\item Перезагружает исходные файлы, изменённые после загрузки
\item Создаёт новый пустой файл программы
\item Собирает проект с помощью утилиты make
\end{itemize}
{\small\textit{Правильный ответ: Перезагружает исходные файлы, изменённые после загрузки.}}

\item Какое сообщение выдаст SWI-Prolog на запрос к неопределённому предикату \texttt{?-} uncle(X, jim).?
\begin{itemize}[label=$\square$, leftmargin=1.5cm, itemsep=0.2em, topsep=0.3em]
\item false.
\item Warning: Singleton variables
\item ERROR: Unknown procedure: uncle/2
\item ERROR: Syntax error: Operator expected
\end{itemize}
{\small\textit{Правильный ответ: ERROR: Unknown procedure: uncle/2.}}

\item О чём предупреждает сообщение Singleton variables: [Y] при загрузке файла?
\begin{itemize}[label=$\square$, leftmargin=1.5cm, itemsep=0.2em, topsep=0.3em]
\item Переменная Y не может быть унифицирована
\item Переменная Y встречается в клаузе только один раз, что часто означает опечатку
\item Переменная Y объявлена глобальной
\item Переменная Y используется в нескольких файлах
\end{itemize}
{\small\textit{Правильный ответ: Переменная Y встречается в клаузе только один раз, что часто означает опечатку.}}

\item Почему правило a(X, Y) :- a(Z, Y), parent(X, Z). приводит к переполнению стека?
\begin{itemize}[label=$\square$, leftmargin=1.5cm, itemsep=0.2em, topsep=0.3em]
\item В теле правила слишком много переменных
\item Это левая рекурсия: рекурсивный вызов стоит первым и выполняется до проверки фактов, поиск уходит в бесконечный спуск
\item Предикат parent/2 не определён
\item SWI-Prolog не поддерживает рекурсию
\end{itemize}
{\small\textit{Правильный ответ: Это левая рекурсия: рекурсивный вызов стоит первым и выполняется до проверки фактов, поиск уходит в бесконечный спуск.}}

\item В каком порядке SWI-Prolog перебирает клаузы предиката и цели в теле правила?
\begin{itemize}[label=$\square$, leftmargin=1.5cm, itemsep=0.2em, topsep=0.3em]
\item Клаузы сверху вниз в порядке записи в файле, цели тела слева направо
\item Клаузы в случайном порядке, цели справа налево
\item Сначала правила, затем факты
\item Клаузы в алфавитном порядке
\end{itemize}
{\small\textit{Правильный ответ: Клаузы сверху вниз в порядке записи в файле, цели тела слева направо.}}

\item Каким запросом можно узнать текущую рабочую директорию в SWI-Prolog?
\begin{itemize}[label=$\square$, leftmargin=1.5cm, itemsep=0.2em, topsep=0.3em]
\item \texttt{?-} current\_dir(D).
\item \texttt{?-} working\_directory(D, D).
\item \texttt{?-} home(D).
\item \texttt{?-} path(D).
\end{itemize}
{\small\textit{Правильный ответ: \texttt{?-} working\_directory(D, D).}}

\item Что вернёт запрос \texttt{?-} findall(X, parent(tom, X), L).?
\begin{itemize}[label=$\square$, leftmargin=1.5cm, itemsep=0.2em, topsep=0.3em]
\item L = bob
\item L = [bob, liz]
\item L = [tom]
\item false.
\end{itemize}
{\small\textit{Правильный ответ: L = [bob, liz].}}

\item Каким запросом открыть встроенную справку по предикату halt/0?
\begin{itemize}[label=$\square$, leftmargin=1.5cm, itemsep=0.2em, topsep=0.3em]
\item \texttt{?-} man(halt).
\item \texttt{?-} info halt.
\item \texttt{?-} help(halt/0).
\item \texttt{?-} doc halt.
\end{itemize}
{\small\textit{Правильный ответ: \texttt{?-} help(halt/0).}}

\item Для чего служит запрос \texttt{?-} apropos(consult).?
\begin{itemize}[label=$\square$, leftmargin=1.5cm, itemsep=0.2em, topsep=0.3em]
\item Загружает файл consult.pl
\item Ищет в документации предикаты и разделы по ключевому слову
\item Выводит исходный код consult/1
\item Отключает предикат consult/1
\end{itemize}
{\small\textit{Правильный ответ: Ищет в документации предикаты и разделы по ключевому слову.}}

\item Какую логическую связку обозначает запятая в теле правила?
\begin{itemize}[label=$\square$, leftmargin=1.5cm, itemsep=0.2em, topsep=0.3em]
\item Конъюнкцию (логическое И)
\item Дизъюнкцию (логическое ИЛИ)
\item Отрицание
\item Импликацию
\end{itemize}
{\small\textit{Правильный ответ: Конъюнкцию (логическое И).}}

\item Какое расширение обычно имеют исходные файлы SWI-Prolog?
\begin{itemize}[label=$\square$, leftmargin=1.5cm, itemsep=0.2em, topsep=0.3em]
\item .swi
\item .pl
\item .prg
\item .lsp
\end{itemize}
{\small\textit{Правильный ответ: .pl.}}

\item Как читается декларативно правило father(X, Y) :- parent(X, Y), male(X).?
\begin{itemize}[label=$\square$, leftmargin=1.5cm, itemsep=0.2em, topsep=0.3em]
\item X — отец Y, если Y — родитель X и Y — мужчина
\item X — отец Y, если X — родитель Y и X — мужчина
\item X — отец Y или X — мужчина
\item Каждый мужчина является отцом
\end{itemize}
{\small\textit{Правильный ответ: X — отец Y, если X — родитель Y и X — мужчина.}}

\item Чем программа на Visual Prolog отличается от программы на SWI-Prolog?
\begin{itemize}[label=$\square$, leftmargin=1.5cm, itemsep=0.2em, topsep=0.3em]
\item В Visual Prolog нет правил, только факты
\item В Visual Prolog обязательны секции domains, predicates и clauses с объявлением типов
\item Visual Prolog не поддерживает рекурсию
\item В Visual Prolog запросы вводятся только в top-level
\end{itemize}
{\small\textit{Правильный ответ: В Visual Prolog обязательны секции domains, predicates и clauses с объявлением типов.}}
\end{enumerate}
\normalsize
